%% ---- Supplementary Note [note:dcell-training] -- OWNED BY its section agent; fill in place. ----
%% Every number below is produced by experiments/006-kuzmin-tmi/scripts/dcell_training_wandb.py
%% from the frozen wandb histories in experiments/006-kuzmin-tmi/results/dcell_training/
%% (runs.csv lists every run id). Figure: FigS-dcell-training.drawio, composed by
%% experiments/006-kuzmin-tmi/scripts/dcell_training_compose_figure.py. Detail cut from the
%% prose on 2026.09.03 (per-stage step times, the dcell_opt rewrite, the per-GPU batch, the
%% partial no-auxiliary run, the epoch-100-to-332 statistics) lives in the figure caption, the
%% two tables, and notes/experiments.006-kuzmin-tmi.scripts.dcell_training_wandb.md.
%% The dataset-build mismatch is stated once, in note:dango-full; this note points there.
\clearpage
\suppnote{note:dcell-training}{DCell training cost, instability, and the checkpoint-based baseline\secstatus{todo}}

The DCell value in Fig.~\ref{fig:ggi}d ($r=0.157\pm0.009$) comes from one training run,
not from replicates: the three values behind the mean and error bar are validation
evaluations of that run at epochs 144, 148, and 150, so the error bar is not a replicate
error bar. DCell in TorchCell (\suppnoteref{note:dcell-model}) took 30.8 days on four GPUs
to reach epoch 333, passed its best validation epoch after 14 days, was unstable
throughout, and was run to completion once within the compute available.

\notesec{Training run and checkpoint statistic} One full run (wandb \texttt{eni948by},
SLURM job 1922684) trained the 2,655-term, 20.6~M-parameter DCell with the DCell loss
(auxiliary losses on, $\alpha=0.3$), AdamW at learning rate $10^{-3}$ and weight decay
$10^{-6}$, global batch 2,400 over four GPUs under DDP, bf16 mixed precision, and gradient
clipping at 10, on the experiment-006 build (\suppnoteref{note:dango-full};
\supptab{tab:dcell-training-cost}). The reduce-on-plateau floor equaled the initial rate,
so the learning rate never changed (\suppfig{fig:dcell-training}b), and the run was stopped
after 333 of the configured 1,000 epochs. The three values in Fig.~\ref{fig:ggi}d, 0.155,
0.142, and 0.173, include the maximum over all 333 per-epoch evaluations
(\supptab{tab:dcell-training-checkpoints}), so their mean (0.157) is an upward-biased order
statistic rather than an expected performance, and their SD (0.016; SEM 0.009) measures
epoch-to-epoch variation of one run, not replicate variance. The checkpoint the training
script saved, the minimum validation MSE at epoch 142, scores $r=0.127$; the final epoch
scores 0.038. After epoch 100, validation~$r$ fluctuated between 0.04 and 0.17 and drifted
down while the training loss kept falling by an order of magnitude
(\suppfig{fig:dcell-training}a,\,b), overfitting under a constant learning rate; whether a
decaying schedule would have stabilized the run is untested.

As for DANGO (\suppnoteref{note:dango-full}), the run on the larger build scored lower: the
one DCell run on the Kuzmin 2018-only build peaked at $r=0.259$ against 0.173 on the
experiment-006 build (\suppfig{fig:dcell-training}f; one run each, differently split, with
auxiliary losses and precision also differing, so the gap is not attributed to data size).

\notesec{Cost and where the time goes} DCell reached its best validation epoch after 1,340
GPU-hours, against 97--101 GPU-hours for each CGT replicate (best at epoch 24--25) and 8--12
for the two DANGO runs on the same build (\supptab{tab:dcell-training-cost};
\suppfig{fig:dcell-training}c): 33 training samples per second (2.22~h per epoch) against 89
for the CGT and about 1,200 for DANGO, and six times as many epochs as the CGT to peak. Three
replicate runs of the same length would cost about 8,900 GPU-hours at the observed rate,
which is why one run is reported and labeled as such. The forward pass visits the 2,655
subsystems one at a time in a Python loop ordered by stratum, so every step launches
hundreds of thousands of small kernels, more as the batch grows; the speed-up work on a shared four-GPU
workstation (\suppfig{fig:dcell-training}d) took a step from 128~s to 20~s, and day-to-day
reruns of one configuration varied by as much as several of the optimizations. A CPU
profile of one step of the same architecture (\suppfig{fig:dcell-training}e) counts 644,603
kernel-level ops at batch 8, about 13,600 more per added strain, because the gene-state
gather runs once per term and per strain; no GPU profile exists. Hypothesis (untested): the
GPU step is bound by those launches rather than by data loading, which the cluster
measurements cannot separate; if so, one grouped matrix multiply per stratum and CUDA
graphs for the fixed per-stratum shapes would remove most of the launch overhead.

%% Figure: notes/assets/drawio/FigS-dcell-training.drawio (composed by
%% experiments/006-kuzmin-tmi/scripts/dcell_training_compose_figure.py from the panels of
%% dcell_training_wandb.py (a-d, f) and dcell_training_cpu_profile.py (e))
%% -> figures/FigS-dcell-training.pdf.
\begin{figure*}[t]
\tcfig{figures/FigS-dcell-training.pdf}{\textbf{DCell training on the trigenic interaction
task: one run, its checkpoint statistic, and its cost.} \textbf{a},~Validation Pearson~$r$ per epoch
for the single full DCell run (wandb \texttt{eni948by}, job 1922684, four GPUs; top axis,
wall-clock) and the stopped run without auxiliary losses (job 1921740, 10 epochs; partial).
Open circles mark the three evaluations averaged in Fig.~\ref{fig:ggi}d; the dotted line is
their mean. \textbf{b},~Training loss (every 10 steps) and validation loss (per epoch) of the same
run; dotted verticals mark the three evaluations; the learning rate was constant.
\textbf{c},~GPU-hours to the best validation epoch (left) and wall-clock hours per epoch (right,
median over epochs) for DCell (four GPUs), DANGO (two runs, two GPUs each) and CGT (three
replicates, four GPUs each); bar, mean; open circles, runs; whisker, SEM where $n>1$
(\supptab{tab:dcell-training-cost}). \textbf{d},~Speed-up stages on a shared four-GPU
workstation: what each stage changed, batch per GPU, samples per second, and seconds per
optimizer step as bars. Arrows mark the cumulative chain 1 to 5, measured in one sitting;
stages 8 and 9 add \texttt{torch.compile} at batch 500 and 600; rows 6 and 7 (stage 5 rerun
on later days, 99 and 119~s) are the dashed range, because the shared machine drifts from
day to day by more than several optimizations and only same-sitting stages compare.
\textbf{e},~Where one training step spends its time on a CPU (Apple M1 Max, fp32, batch 8), a
stand-in until the pending GPU profile on gilahyper replaces it. The step issues 644,603
leaf ops (67 distinct), each a kernel launch on a GPU, growing by about 13,600 per added
strain; the per-term gene-state gather and 2,655 BatchNorm calls cost more than the matrix
multiplies. \textbf{f},~Best validation Pearson~$r$ (maximum over epochs) of the same DCell on
the Kuzmin 2018-only build (91,050 records; \texttt{biucpv7p}, best at epoch 227) and the
experiment-006 build (332,313; \texttt{eni948by}, best at epoch 150), one run each, splits
and training settings differing; the only measured data-size comparison for DCell.}{fig:dcell-training}
\end{figure*}

%% SOURCE: experiments/006-kuzmin-tmi/scripts/dcell_training_wandb.py -- AUTO-GENERATED, do not hand-edit; rerun the script.
\begin{table}[t]
\centering
\footnotesize
\caption{The DCell value in Fig.~\ref{fig:ggi}d and the run it comes from. All values are
the validation Pearson~$r$ logged by the single DCell training run (wandb
\texttt{eni948by}, SLURM job 1922684, 333 epochs, 30.8 days on 4 GPUs). \emph{Rank} is the value's rank among
the 333 per-epoch validation evaluations of that run. The three epochs are
evaluations of one run, so their spread is checkpoint-to-checkpoint variation, not replicate
variance, and the reported SEM is not a replicate SEM.}
\label{tab:dcell-training-checkpoints}
\begin{tabular}{@{}l r r r r@{}}
\toprule
\textbf{Evaluation} & \textbf{Epoch} & \textbf{Wall-clock (d)} & \textbf{Val.\ Pearson $r$} & \textbf{Rank}\\
\midrule
Fig.~\ref{fig:ggi}d checkpoint & 144 & 13.4 & 0.155 & 2\\
Fig.~\ref{fig:ggi}d checkpoint & 148 & 13.8 & 0.142 & 4\\
Fig.~\ref{fig:ggi}d checkpoint & 150 & 14.0 & 0.173 & 1\\
\midrule
Mean of the three (Fig.~\ref{fig:ggi}d) & & & 0.157 & \\
SD over the three (ddof$=$1) & & & 0.016 & \\
SEM $=$ SD$/\sqrt{3}$ (Fig.~\ref{fig:ggi}d error bar) & & & 0.0091 & \\
\midrule
Maximum over all epochs & 150 & 14.0 & 0.173 & 1\\
Minimum validation MSE (the saved checkpoint) & 142 & 13.2 & 0.127 & 20\\
Final epoch & 332 & 30.8 & 0.038 & 231\\
\bottomrule
\end{tabular}
\end{table}

%% SOURCE: experiments/006-kuzmin-tmi/scripts/dcell_training_wandb.py -- AUTO-GENERATED, do not hand-edit; rerun the script.
\begin{table*}[t]
\centering
\footnotesize
\caption{Training cost of the three learned models compared in Fig.~\ref{fig:ggi}d, from the
wandb histories of the runs listed in \texttt{results/dcell\_training/runs.csv}. Samples per
epoch is optimizer steps per epoch times the global batch. \emph{Best} is the epoch of the
maximum validation Pearson~$r$ over the run; the wall-clock and GPU-hours to reach it are the
logged run time at that epoch. CGT rows are the three replicate training jobs whose
best-Pearson checkpoints give the Fig.~\ref{fig:ggi}d value; DANGO rows are the two 1{,}000-epoch training
runs on the same build as DCell. DCell and DANGO trained on the experiment-006 build
(332,313 records); the CGT replicates trained on the experiment-010 build (376,732 records).}
\label{tab:dcell-training-cost}
\begin{tabular}{@{}l l l r r r r r r r r@{}}
\toprule
\textbf{Model} & \textbf{Run} & \textbf{Build} & \textbf{GPUs} & \textbf{Samples/epoch} & \textbf{h/epoch} & \textbf{Samples/s} & \textbf{Best epoch} & \textbf{Best $r$} & \textbf{h to best} & \textbf{GPU-h to best}\\
\midrule
DCell & \texttt{eni948by} & 006 & 4 & 266,393 & 2.22 & 33 & 150 & 0.173 & 334.9 & 1,340\\
CGT & \texttt{lzs9pcj3} & 010 & 4 & 302,064 & 0.94 & 90 & 24 & 0.452 & 24.3 & 97\\
CGT & \texttt{yv4r30bi} & 010 & 4 & 302,064 & 0.94 & 89 & 25 & 0.447 & 25.3 & 101\\
CGT & \texttt{c7671wgj} & 010 & 4 & 302,059 & 0.94 & 89 & 24 & 0.462 & 24.5 & 98\\
DANGO & \texttt{014mprap} & 006 & 2 & 265,856 & 0.06 & 1,244 & 98 & 0.368 & 6.0 & 12\\
DANGO & \texttt{9jpfy547} & 006 & 2 & 265,856 & 0.07 & 1,134 & 60 & 0.354 & 4.1 & 8\\
\bottomrule
\end{tabular}
\end{table*}

